% ============================================================================
% APÊNDICE F — Guia de revisão e contribuição
% ============================================================================
\chapter{Revisão por Pares e Contribuição}
\label{apendice:revisao}

\begin{caixaconceito}
Este livro é revisado por um \textbf{sistema de revisão por pares simulado}:
agentes independentes atuam como revisores técnico, pedagógico e de fontes,
aplicando checklists objetivos e validando o resultado (compilação LaTeX,
execução dos códigos, cruzamento das referências). O processo completo está
documentado em \texttt{revisao/README.md}.
\end{caixaconceito}

\section{As três revisões}

\begin{enumerate}[leftmargin=*, itemsep=0.4em]
  \item \textbf{Revisão técnica}: os exemplos funcionam? A versão citada está
  correta? As APIs usadas existem na versão documentada? O código executa e os
  testes passam?
  \item \textbf{Revisão pedagógica}: a ordem é linear e progressiva? Os
  pré-requisitos estão respeitados? Exercícios têm dificuldade crescente?
  Analogias e exemplos são claros?
  \item \textbf{Revisão de fontes}: cada dado de mercado tem fonte oficial?
  Cada afirmação técnica remete à documentação? As datas de consulta estão
  registradas?
\end{enumerate}

\section{O que a revisão valida (checklist resumido)}

\begin{itemize}[leftmargin=*, itemsep=0.4em]
  \item Compilação LaTeX sem erros e sem \emph{overfull hbox} significativo;
  \item \texttt{tsc --noEmit} e testes (Vitest/Playwright) verdes no repositório;
  \item Todas as referências cruzadas (\texttt{\textbackslash ref}) resolvem;
  \item Nenhuma caixa didática quebrada ou lista estourada;
  \item Nenhuma afirmação sem fonte nas seções de mercado;
  \item Exercícios com solução disponível e testada;
  \item Termos técnicos em inglês com explicação em pt-BR consistente.
\end{itemize}

\section{Como contribuir}

\begin{itemize}[leftmargin=*, itemsep=0.4em]
  \item \textbf{Erros de código}: abra uma issue com o capítulo, a listagem e
  o erro reproduzido;
  \item \textbf{Atualizações de versão}: PR alterando o exemplo + atualizando
  a tabela de versões (apêndice~\ref{apendice:manutencao});
  \item \textbf{Melhorias didáticas}: sugira analogias, exercícios e projetos;
  \item \textbf{Tradução/adaptação}: o livro é CC BY 4.0 — adapte citando a fonte.
\end{itemize}

\begin{caixadica}
O repositório \texttt{revisao/} contém os checklists completos, os scripts de
validação (\texttt{scripts/}) e um modelo de relatório de revisão
(\texttt{revisao/relatorio-modelo.md}) para registrar cada rodada.
\end{caixadica}
